\begin{abstract}
Electricity-market data are graded and released field by field, yet the inference risk of a field
depends on what else an adversary knows: fields that reveal nothing alone can, together with other
released data, expose unit outputs, private offers, line flows or load before release. We define
field-level risk from set-level inferability, the out-of-sample accuracy that attackers retrained on
a set of fields reach for a sensitive target. The contextual gain of a field over a background has a
whole distribution; attribution methods average it, whereas a security audit needs its worst case.
The budgeted marginal risk $M^{(K)}$ is the largest gain over any background of at most $K$ fields,
relative to data already public, and its excess over single-field risk is the background
amplification $\Gamma^{(K)}$. $M^{(K)}$ is a budget-constrained marginal contribution index, the
minimal threshold-free safety margin, and flags exactly the members of minimal unsafe combinations of
at most $K+1$ fields; its top backgrounds tell the data owner what not to co-release. To evaluate all
small sets repeatedly, an amortized attacker---one network pre-trained under random visibility masks,
queried through a closed-form readout per field set---scans the sets, and retraining a few backgrounds
per field certifies the result. In a controlled market case where one unit's private offer markup is
the target, every price field has zero risk alone and all of its risk comes from the unit's output as
background, ranked and regime-dependent as the dispatch physics implies; correlation, single-field and
attribution scores miss it. On four power data sets with rule-based field roles and
$1.2\times10^5$ exhaustively retrained field sets as ground truth, amplification is common---a typical
field's average contextual gain is a third of its worst case---and single-field grading recovers only
13--17\% of the critical fields, whereas scan--certify recovers up to 91\% without false positives.
We also find that masked pre-training on the targets alone collapses the representation and, on the
small sets of an audit, adds little over random features; masked reconstruction of all columns, features concatenated with
those of an untrained network and a safeguard against exploding predictions reduce the error of field risk by 29\% relative to target-only pre-training.
\end{abstract}
